
\chapter{Transformation Audit M001}
\label{chap:audit-m001}

\section{Mathematical Specification}

\subsection*{Scientific Purpose}
Construct an orientation-preserving integer basis adapted to a primitive
Miller plane.

\subsection*{Domain}
Primitive Miller indices $\mathbf h\in\mathbb Z^3$ with
$\gcd(h,k,l)=1$.

\subsection*{Codomain}
Matrices $S\in SL(3,\mathbb Z)$ satisfying
\[
\mathbf hS=(0,0,1).
\]

\subsection*{Scientific Invariants}
\begin{itemize}
\item Integer lattice preserved.
\item Orientation preserved.
\item Surface normal preserved.
\item Kernel basis remains primitive.
\end{itemize}

\subsection*{Gauge Freedom}
Changing the sign of an in-plane basis vector does not alter the surface
lattice and is treated as gauge fixing.

\section{Algorithm}

\begin{enumerate}
\item Reduce the Miller index to primitive form.
\item Construct Bézout coefficients.
\item Build a unimodular reduction.
\item Permute the basis to align the surface normal.
\item Repair orientation if necessary.
\item Verify all defining invariants.
\end{enumerate}

\section{Implementation Mapping}

\begin{center}
\begin{tabular}{lll}
\hline
Algorithm step & CALM procedure & Notes\\
\hline
Primitive reduction & \texttt{\_reduce\_hkl()} & Exact integer arithmetic\\
Bézout coefficients & \texttt{\_extended\_gcd()} & Extended Euclidean algorithm\\
Unimodular reduction & \texttt{\_right\_reduce\_row\_to\_hnf()} & Successive Bézout elimination\\
Assembly & \texttt{surface\_basis\_S\_from\_hkl()} & Composite transformation\\
\hline
\end{tabular}
\end{center}

\section{Verification Matrix}

\begin{center}
\begin{tabular}{lll}
\hline
Obligation & Evidence & Status\\
\hline
Primitive Miller index & Implementation & Review complete\\
$\mathbf hS=(0,0,1)$ & Assertions/tests & Verified\\
$\det(S)=1$ & Implementation & Verified\\
Primitive kernel basis & Requires dedicated audit & Pending\\
$S^{-1}\in SL(3,\mathbb Z)$ & Not yet identified & Pending\\
\hline
\end{tabular}
\end{center}

\section{Audit Findings}

\begin{enumerate}
\item The implementation decomposes naturally into four mathematical
sub-transformations rather than a single operation.
\item Orientation repair is a gauge-fixing operation, not a scientific
morphism.
\item The routine \texttt{\_right\_reduce\_row\_to\_hnf()} is best
understood as a constructive unimodular reduction based on successive
Bézout eliminations.
\end{enumerate}
